% Proofs added in the unification revision: answers over objects, loading (Propositions 5.1-5.3,
% Corollary 5.4 of the paper) and the one-semantics theorem for ORMatic's translation (Theorem 6.1).

% ============================================================
\section{Answers over Objects and over Individuals}
\label{sec:objects-individuals}

EQL evaluates $=$, $\neq$ and $\in$ with Python's \texttt{operator.eq} and
\texttt{operator.contains}, so the classes of the compared objects decide equality. We therefore
define $\mathit{ans}(Q,\mathcal{K})$ over assignments $\theta$ of \emph{objects} and values, with
$o=o'$ read as the equality of their classes. In the current version, roles and classes declared with \texttt{eq=False} compare by identity
(\texttt{Role.\_\_eq\_\_} is \texttt{self is other}, \texttt{patterns/role.py}), while a plain
dataclass compares its fields; the predicate \texttt{IsSameSemanticEntity}
(\texttt{patterns/role\_predicates.py}) compares the role takers of two objects. In the
earlier version, which ran the experiments, a role's hash is the identity of its root taker and
its equality compares hashes (\texttt{class\_diagrams/utils.py}, template \texttt{onto\_base.j2}),
so a role equals its taker and its taker's other roles. In both versions two root objects that
share one IRI are unequal.

\begin{definition}
The \emph{answers over individuals} of $Q$ are
$\mathit{ind}(\mathit{ans}(Q,\mathcal{K}))=\{\mathit{ind}\circ t \mid t\in\mathit{ans}(Q,\mathcal{K})\}$.
\end{definition}

\begin{proposition}
\label{prop:objects-individuals}
Let $V_\alpha$ be the set of values the two operands of an atom $\alpha$ of the form $t_1=t_2$,
$t_1\neq t_2$ or $t_2\in t_1$ can take. If, for every such atom of $Q$, $\mathit{ind}$ is
one-to-one on $V_\alpha$ and object equality on $V_\alpha$ is identity or coincides with equality
of $\mathit{ind}$, then $\mathit{ind}(\mathit{ans}(Q,\mathcal{K}))$ equals the answers of $Q$
evaluated over individuals in $\mathcal{I}_{\mathcal{A}}$.
\end{proposition}

\begin{proof}
By induction on $\varphi$. Class atoms, property atoms and predicates are defined through
$\mathit{ind}$ (Section~\ref{sec:eql:semantics}), so their truth under $\theta$ and under
$\mathit{ind}\circ\theta$ agrees. For $t_1=t_2$ with values $u,v\in V_\alpha$: $u=v$ as objects
iff $u$ is $v$ (or $\mathit{ind}(u)=\mathit{ind}(v)$), iff $\mathit{ind}(u)=\mathit{ind}(v)$, since
$\mathit{ind}$ is one-to-one on $V_\alpha$; $\neq$ and $\in$ follow, the latter as a disjunction
of equalities over the collection's elements. Connectives and quantifiers commute with the
induction hypothesis, because every individual a quantifier ranges over is the image of an object
of the quantified domain (domain closure).
\end{proof}

The hypothesis holds, for example, when every individual is denoted by one object plus its roles,
at most one of each class, and both operands of each atom range over one class. Positive queries
contain no $\neq$ and no negation; for them, object equality implies equality of the denoted
individuals in both versions, so every answer over objects maps to an answer over individuals
without the hypothesis.

% ============================================================
\section{Loading: Proofs of Propositions 5.1--5.3 and Corollary 5.4}
\label{sec:loading-proofs}

We write $F_k$ for the facts (class and property assertions over individuals) after step $k$ of
Algorithm~1 of the paper, $R$ for the assign rules (\texttt{prp-spo1}, \texttt{prp-spo2},
\texttt{prp-eqp1/2}, \texttt{prp-inv1/2}, \texttt{prp-symp}, \texttt{prp-trp}), and
$\mathit{Imp}^*(P)$ for the properties reachable from $P$ through super-, equivalent and inverse
properties and chains. The constructs covered are: named classes with subclass and equivalence
axioms; domains and ranges that are named classes; sub-, equivalent, symmetric and transitive
properties, one inverse per property, and property chains; restrictions and intersections as
superclasses; unions in subclass position; and one sufficient condition per class, a
\texttt{someValuesFrom} or \texttt{hasValue} restriction with a named filler, possibly intersected
with named classes. Of the 78 RL/RDF rules (Section~\ref{sec:rl-rules}), 14 apply when loading, 9
on assignment, 15 by construction of the hierarchies, 21 are checked, 13 derive nothing new here,
and 6 concern constructs outside this list.

\paragraph{Hypotheses of Proposition 5.2} (G1) every property lists all its equivalent properties
(the generator lists the declared ones, and step~4 does not apply equivalence to facts derived by
equivalence); (G2) equivalent properties have the same characteristics; (G3) no symmetric,
transitive property is implied, through a chain, by another symmetric, transitive property
(step~5 is a single pass); (G4) every object IRI is declared a named individual; every domain and
range is a named class; every individual has at most one value of each data property. OWL2Bench
meets G1--G4.

\begin{lemma}[Steps 1 and 2]
\label{lem:step12}
$F_1$ contains, for every asserted $P(a,b)$, the facts that \texttt{prp-spo1}, \texttt{prp-eqp1/2},
\texttt{prp-inv1/2} and \texttt{prp-symp} derive from it, closed under these rules. $F_2$ is
additionally closed under \texttt{prp-spo2} and under \texttt{prp-trp} for every transitive
property in $\mathit{Imp}^*$ of a property used by a class expression or a chain, and under the
typing rules \texttt{prp-dom}, \texttt{prp-rng}, \texttt{cax-sco}, \texttt{cax-eqc1/2} and the
\texttt{cls-*} rules of the covered constructs.
\end{lemma}

\begin{proof}
Step~1 builds each proxy from the asserted types and values and those implied by sub-, inverse
and symmetric properties (\texttt{owl\_instances\_loader.py}); equivalence is applied with G1.
Step~2 repeats, until no type or value is added, the chain and transitivity closure of the
relevant properties and the typing rules; each iteration is monotone over a finite set of
facts, so it reaches the least fixpoint, which is closed under the rules it applies.
\end{proof}

\begin{lemma}[Typing invariant]
\label{lem:typing}
Every fact in the final $\mathcal{A}^*$ that is not in $F_2$ is a fact $P(a,b)$ of a property $P$
in $\mathit{Imp}^*(T)$ for a transitive property $T$ that no class expression or chain uses, and
$a$ and $b$ already have facts of $P$, as subject and as object respectively, in $F_2$. Hence the
typing rules derive no fact that is not in $F_2$.
\end{lemma}

\begin{proof}
Facts added in steps~4 and~5 come from $R$ applied to facts of $F_2$ (step~4) or from component
pairs of a symmetric, transitive property (step~5). \texttt{prp-spo1}, \texttt{prp-eqp},
\texttt{prp-inv} and \texttt{prp-symp} were already closed in $F_1$ (Lemma~\ref{lem:step12}),
so a new fact arises only from \texttt{prp-trp} (or a chain over its result, closed in step~2 for
the used properties) on a non-relevant transitive $T$, and $T(a,c)$ from $T(a,b),T(b,c)$ relates
a subject and an object that already hold $T$; its super-, equivalent and inverse facts relate
individuals that already hold those properties, by the closure of $F_1$. Step~5 relates members
of one weakly connected component, each of which already has a fact of the property in both
positions (symmetry). A typing rule's premise $P(a,b)$ therefore has a counterpart $P(a,b')$ and
$P(a'',b)$ in $F_2$, whose conclusions are already in $F_2$; \texttt{hasValue} conclusions are
facts of $F_2$'s types, added after step~4.
\end{proof}

\begin{lemma}[Storage]
\label{lem:storage}
Every fact of $F_2$ and every fact that steps~4 and~5 derive is stored: its subject has an object
whose class declares the property.
\end{lemma}

\begin{proof}
TBox compilation declares each property on its domain class, and on the domains derived through
equivalent, inverse and super-properties (\texttt{prp-dom}, \texttt{scm-dom2}), which are entailed
by the ontology. By Lemma~\ref{lem:step12}, the subject of every fact of $F_2$ has that domain
among its types, so step~3 creates an object of a class that declares the property; by
Lemma~\ref{lem:typing} the same holds for later facts.
\end{proof}

\begin{lemma}[Step 4]
\label{lem:step4}
After step~4, the stored property facts are closed under $R$, except \texttt{prp-trp} for
symmetric, transitive properties.
\end{lemma}

\begin{proof}
Every fact is added to the store before it is enqueued, and every enqueued fact is propagated
(Definition~\ref{def:assign}), so for every instance of a rule of $R$ whose premises are present
at the end, consider the premise added last: when it was propagated, all other premises were
present, so the conclusion was added. The implementation skips an instance only when its
conclusion is certainly present: the symmetric fact of a fact derived by symmetry, the inverse
of a fact derived by inverse (one inverse per property), and the equivalence of a fact derived by
equivalence (all equivalents listed, G1). The one remaining skip, transitivity on a fact derived
from an equivalent property, is covered by G2: the equivalent twin has the same characteristics,
and its own propagation applies transitivity to the same pair.
\end{proof}

\begin{lemma}[Step 5]
\label{lem:step5}
For a symmetric, transitive $P$, the closure of $P$'s facts under \texttt{prp-symp} and
\texttt{prp-trp} relates exactly the pairs of members of each weakly connected component of $P$'s
graph; with G3, step~5 completes all of them in one pass, and propagates the new facts through
$R$ when $P$ has a super-property, an inverse, an equivalent or occurs in a chain.
\end{lemma}

\begin{proof}
A path in the underlying undirected graph between $a$ and $b$ yields $P(a,b)$ by symmetry and
transitivity, and conversely every derived pair lies on such a path. G3 ensures that no other
symmetric, transitive property gains facts from the propagation, which a single pass would miss.
\end{proof}

\begin{proof}[Proof of Proposition 5.1]
Every fact added by steps~1--5 and by an assignment is the conclusion of a sequence of RL/RDF
rule applications (Lemmas~\ref{lem:step12}--\ref{lem:step5}; the schema-level domains and ranges
are entailed, Lemma~\ref{lem:storage}), so it is entailed by $\mathcal{O}$. A positive query is
preserved under homomorphisms, and $\mathcal{I}_{\mathcal{A}^*}$ maps homomorphically into every
model of $\mathcal{O}$, also one that identifies individuals KRROOD keeps apart; with
Section~\ref{sec:objects-individuals}, the image under $\mathit{ind}$ of an answer is a certain
answer.
\end{proof}

\begin{proof}[Proof of Proposition 5.2]
By Theorem~PR1 of the OWL~2 profiles, it suffices that $\mathcal{A}^*$ is closed under the
RL/RDF rules for the covered constructs. The check reports no application of an equality rule,
so these rules derive nothing. The property rules are closed by Lemmas~\ref{lem:step4}
and~\ref{lem:step5}, the typing rules by Lemmas~\ref{lem:step12} and~\ref{lem:typing}, and every
fact is stored by Lemma~\ref{lem:storage}. G4 makes every individual an object and every data
value a field.
\end{proof}

\begin{proof}[Proof of Proposition 5.3]
An assignment propagates the new fact as in Lemma~\ref{lem:step4}, from the current facts; the
same argument (premise added last, G2) gives closure under $R$, except transitivity of
symmetric, transitive properties, which only step~5 completes, and except facts whose subject has
no object declaring the property, which are not stored. Soundness follows as in
Proposition~5.1, with the asserted facts added to $\mathcal{O}$.
\end{proof}

\begin{proof}[Proof of Corollary 5.4]
The earlier loader may store, as a value, a role whose class fits the range through its taker, so
a stored value's class can differ from the declared range; the earlier version's equality, which
equates a role with its taker and sibling roles, makes such values equal to the variable's object
of the same individual.
By Proposition~5.2, every certain answer $\bar a$ of a positive query $Q$ satisfies $Q$ over
individuals in $\mathcal{I}_{\mathcal{A}^*}$. Under the corollary's hypothesis,
Proposition~\ref{prop:objects-individuals} applies, so $\bar a$ is the image of an answer over
objects.
\end{proof}

% ============================================================
\section{One Semantics for Working and Long-Term Memory}
\label{sec:ormatic-theorem}

\paragraph{The database $D(\mathcal{A})$} ORMatic's \texttt{to\_dao} converts all objects of
$\mathcal{A}$, with one shared conversion state, into an empty database: joined-table
inheritance with one table per class, an association table per collection attribute, a nullable
foreign key per reference, a role table with the column \texttt{\_role\_taker\_id}, the IRI in
the \texttt{uri} column of \texttt{ThingDAO}, which is \texttt{NOT NULL}, and inferred facts stored
like asserted ones.

\paragraph{The fragment $\mathcal{F}$} Queries \texttt{an}/\texttt{the} over \texttt{entity} or
\texttt{set\_of} of variables, collection elements (\texttt{flat\_variable}) and data attributes,
with an optional plain \texttt{.distinct()}. Terms: reference chains, including delegation from a
role to its taker and narrowing to the single subclass that declares an attribute. Atoms:
$=,\neq,<,\le,>,\ge$ between a column and a literal, two columns, or two elements of one DAO
hierarchy; \texttt{contains} on an object collection; \texttt{in\_} with a literal list; an
attribute as a truth value. Connectives: $\wedge$; $\vee$ and \texttt{not\_} without joins inside
them (the translator's guard), negation translated as $(\tau\varphi)$ \texttt{IS NOT TRUE};
\texttt{exists} over an element chain or a variable, as a correlated \texttt{EXISTS}.

\paragraph{Left out, and why} The translator rejects universal quantification (no
translation implemented), predicates and symbolic functions (they run Python code), rules (they
create objects), grouping and aggregates (no proof given here), and, since this revision, a
variable used as a condition (it has no column to test). It accepts but the theorem excludes:
substring \texttt{contains}, which becomes \texttt{LIKE}, where \texttt{\%} and \texttt{\_} are
wildcards and SQLite compares case-insensitively; \texttt{distinct(on=\dots)}, whose \texttt{on}
is ignored; \texttt{limit}, which picks a subset that depends on the evaluation order; and
\texttt{ordered\_by} without \texttt{limit} does not change the set of answers, so it is
admitted as a no-op on answer sets.

\paragraph{Side conditions} (1) The database is $D(\mathcal{A})$. (2) Every variable ranges over
all instances of its class: the translator ignores an explicit \texttt{domain=} and ranges over
all rows of the class's table. (3) Stored classes use single inheritance (the translator joins
one parent table). (4) No missing value (\texttt{None}/\texttt{NULL}) in a compared column or an
\texttt{in\_} list, no \texttt{None} element of a collection, and every role has a taker.
(5) Order comparisons only on numbers, or a database with C collation, and no NaN. (6) Compared
objects use identity equality (\texttt{eq=False}). (7) Answers are compared as sets of tuples of
IRIs and values.

\begin{theorem}[One semantics, Theorem 6.1 of the paper]
For every $Q\in\mathcal{F}$ and every $\mathcal{A}$ meeting (1)--(7), the identifier-selecting
translation of $Q$, evaluated over $D(\mathcal{A})$, returns
$\{\iota\circ t \mid t\in\mathit{ans}(Q,\mathcal{K})\}$, where $\iota$ maps an object to its IRI
and a value to itself.
\end{theorem}

\begin{proof}
By structural induction on $Q$, with the invariant: for a formula $\varphi$ whose free variables
and elements are $\bar z$, the rows of the outer \texttt{FROM} clause correspond one to one to
the assignments of $\bar z$ to objects (and elements to pairs of an owner and a member), and the
SQL condition $\tau\varphi$ is true on a row iff $\varphi$ holds under the corresponding
assignment.

\emph{Base.} By (1) and (2), the rows of a class's table, joined with its parent tables by the
inheritance condition, correspond one to one to the instances of the class; by (3) one join chain
suffices. An element of a collection corresponds to a row of its association table (one row per
pair, as collections are sets). A reference chain becomes joins on foreign keys; an
\texttt{IS NOT NULL} test or the narrowing join to the declaring subclass removes exactly the rows
whose object lacks the attribute, as EQL drops a binding when an attribute is missing; role
delegation joins \texttt{\_role\_taker\_id}, total by (4). A comparison of two columns or a
column and a literal has the same truth value in SQL and Python by (4) (no \texttt{NULL}), (5)
(the same order) and (6) (object comparison is identity, hence comparison of primary keys);
\texttt{in\_} with a literal list is a disjunction of such comparisons. \texttt{contains} on an
object collection is a correlated \texttt{EXISTS} over the association table, true iff the pair
is a row, iff the member is in the set; as a top-level conjunct with \texttt{distinct}, it may be
a join, which yields the same set of answer rows after duplicate removal (7). An attribute as a
truth value is \texttt{IS NOT NULL} for a reference, a non-empty \texttt{EXISTS} for a
collection and the column's truth for a value, as Python's truth value of each.

\emph{Step.} $\wedge$ is \texttt{AND} over the same rows. For $\vee$ and \texttt{not\_}, the
translator's guard ensures that no join occurs inside, so the rows are those of the enclosing
formula and \texttt{OR} and $(\tau\varphi)$ \texttt{IS NOT TRUE} compute the truth value: the
latter is true exactly when $\varphi$ fails, also when SQL evaluates $\tau\varphi$ to
\texttt{NULL}, matching negation as failure. \texttt{exists} over a variable or element chain is
a correlated semi-join, true iff some row of the inner \texttt{FROM} clause satisfies the inner
condition, iff some assignment of the bound variable does. Finally, the selection maps each row
to the IRIs of the selected objects (by (1), the \texttt{uri} column) and the selected columns,
and by (7) duplicates do not matter.
\end{proof}

The check of the 18 OWL2Bench queries (\texttt{run\_ormatic\_translation.py --selection iris})
confirms the theorem empirically: it evaluates each query object in working memory and, through
\texttt{eql\_to\_sql}, in PostgreSQL over the objects built from the closure, with every domain
equal to all instances of its class, and both return GraphDB's answers.

\paragraph{Fixed in this revision} An attribute typed with its own class (e.g.\
\texttt{Node.parent: Node}) was generated without SQLAlchemy's \texttt{remote\_side}: after
\texttt{a.parent = b} and \texttt{c.parent = b}, the database held \texttt{b}$\to$\texttt{a} and
lost \texttt{c}'s fact, which violated (1). ORMatic now marks the referenced primary key, and the
test \texttt{test\_references\_to\_one\_object\_of\_the\_same\_class\_are\_kept} covers it.
